\section{Forelesning 9: \emph{Oppsummering og eksamensforberedelse}}
\label{sec:lec09}
\settermlabel{terms:lec09}

\subsection*{Oversikt}
Den avsluttende forelesningen (Ulrik Røhl \& Sebastian Winslow, 21. april) gjør
en integrativ pass over hele kurset. I stedet for å introdusere nytt
materiale, kobler den de foregående forelesningene sammen til et samlet
analytisk apparat for IT-strategi. Forelesningen inneholder også konkrete
råd til eksamen, hentet fra erfaringene med case competition.

\subsection{Hovedteori -- den integrative tråden}

\subsubsection*{Den røde tråden i kurset}
Kurset bygger en kjede av begreper:
\begin{enumerate}
  \item \autoref{sec:lec01}: \emph{Hva er strategi?} -- bevisst posisjonering
    (Porter 1996), 5 P-er (Mintzberg), bred vs.\ smal digital strategi
    (Chen, Bharadwaj).
  \item \autoref{sec:lec02}: \emph{Hvordan analysere?} -- PESTEL, Five Forces,
    livssyklus, verdikjede, VRIO, SWOT, og katalog av strategiske virkemidler.
  \item \autoref{sec:lec03}: \emph{Hva er forretningsmodellen?} -- Johnson, Weill
    \& Woerners fire DBM-typer, BMC, DVC.
  \item \autoref{sec:lec04}: \emph{Hvilken type endring?} -- Wessel et al.\
    sin distinksjon DT vs.\ ITEOT; Kemps situering av AI.
  \item \autoref{sec:lec05}: \emph{Innovasjon eller disrupsjon?} -- digital
    innovasjon som rekombinasjon; Christensen om disrupsjon; DgI/DgD/DI.
  \item \autoref{sec:lec06}: \emph{Verdikjede eller økosystem?} --
    nettverkseffekter, plattformstyring, three-dimensional chess, PBMC.
  \item \autoref{sec:lec07}: \emph{Hvordan implementere?} -- Tawse \& Tabesh,
    Miller, Kotter, Schein, governance, legacy debt.
\end{enumerate}

\keypoint{Sammen utgjør dette en pipeline: \emph{posisjon $\rightarrow$
analyse $\rightarrow$ forretningsmodell $\rightarrow$ transformasjonstype
$\rightarrow$ innovasjon/disrupsjon $\rightarrow$ økosystem-rolle $\rightarrow$
implementering}. Hvert ledd kan stilles spørsmål om.}

\subsubsection*{Hvordan rammeverkene fungerer sammen}
\begin{itemize}
  \item \emph{Porter} forteller hva strategi \emph{er} (posisjonering, valg,
    fit).
  \item \emph{Weill \& Woerner} forteller hvor virksomheten \emph{står}
    (DBM-type, customer knowledge).
  \item \emph{DVC} forteller hvor digital verdi \emph{kan oppstå}.
  \item \emph{Wessel et al.} forteller hvilken \emph{type endring} digital
    teknologi medfører (DT/ITEOT).
  \item \emph{Yoo et al.\ / Hukal \& Henfridsson} forteller hvordan
    innovasjon \emph{skjer} (rekombinasjon).
  \item \emph{Christensen} forteller når og hvorfor etablerte aktører
    \emph{feiler}.
  \item \emph{Økosystemlitteraturen} (Parker, Eisenmann) forteller hvilken
    \emph{rolle} man spiller (driver, modular producer, supplier).
  \item \emph{Tawse \& Tabesh}, \emph{Miller}, \emph{Kotter},
    \emph{Bjørn-Andersen \& Hedman} forteller hvordan strategi
    \emph{realiseres}.
\end{itemize}

\subsubsection*{Hva rammeverk \emph{kan} og \emph{ikke kan}}
\keypoint{Teorier gir analytiske linser -- de strukturerer \emph{hva man ser
etter}. De kan ikke bevise empiriske påstander om en spesifikk virksomhet.}
Begrensningene:
\begin{itemize}
  \item Teorier er forenklinger.
  \item Virkelige virksomheter passer ikke pent i én boks.
  \item Et rammeverk fra én kontekst overføres ikke perfekt til en annen.
  \item Evidensen må alltid hentes inn separat -- og kildekritikken må
    følge med.
\end{itemize}

\subsection{Sentrale fagbegreper -- konsolidering}\label{terms:lec09}
\begin{itemize}
  \item \term{Strategic positioning}, \term{trade-offs}, \term{fit}, \term{5 Ps}.
  \item \term{PESTEL}, \term{Five Forces}, \term{value chain}, \term{VRIO},
    \term{SWOT}.
  \item \term{Supplier}, \term{Omnichannel}, \term{Modular producer},
    \term{Ecosystem driver} (Weill \& Woerner).
  \item \term{Value proposition}, \term{value creation}, \term{value capture}.
  \item \term{Digital Transformation} vs.\ \term{IT-enabled transformation}
    (Wessel et al.); \term{situating activities} (Kemp).
  \item \term{Digital innovation}, \term{recombination}, \term{disruptive /
    sustaining innovation}, DgI/DgD/DI.
  \item \term{Network effects} (same-/cross-side, +/--), \term{platform
    governance}, \term{boundary resources}, \term{platform envelopment}.
  \item \term{IT governance}, \term{hybrid governance}, \term{legacy debt},
    \term{change management}, \term{Tawse \& Tabesh conditions/actions/
    capabilities}.
\end{itemize}

\subsection{Coop som integrerende eksempel}

Coop fungerer som et gjennomgående case:
\begin{itemize}
  \item \textbf{Posisjon (\autoref{sec:lec01})}: store-first omnichannel --
    et bevisst valg, bekreftet ved lukking av Coop.dk.
  \item \textbf{Analyse (\autoref{sec:lec02})}: høy bransjerivalisering,
    sterk leverandørmakt, økende digital substitusjonstrussel, moden
    bransjefase. Medlemsdata som VRIO-kandidat.
  \item \textbf{Forretningsmodell (\autoref{sec:lec03})}: klassisk
    omnichannel; Lobyco som potensiell modular producer; BMC + DVC som
    analytiske verktøy.
  \item \textbf{Transformasjonstype (\autoref{sec:lec04})}: IT-enabled
    transformation -- ikke digital transformation.
  \item \textbf{Innovasjon (\autoref{sec:lec05})}: Coop App er sustaining +
    digital innovasjon (rekombinasjon), ikke disruptiv.
  \item \textbf{Økosystem (\autoref{sec:lec06})}: ikke en ecosystem driver
    i dag; Lobyco som mulig modular producer i andres økosystem.
  \item \textbf{Implementering (\autoref{sec:lec07})}: legacy-kompleksitet
    (Bjørn-Andersen \& Hedman 2016); hybrid governance og pilotering som
    anbefalt implementeringsdisiplin.
\end{itemize}

\subsection{Eksamenstips fra Lektion 9}
Fra case competition-erfaringene:
\begin{enumerate}
  \item \textbf{Les og svar på hele oppgaven} -- \guillemotleft{}gå inn i casen\guillemotright{}, ikke
    bare reproduser teorien.
  \item \textbf{Vær kritisk til ditt eget svar} -- still motspørsmål.
  \item \textbf{Vær tydelig og konkret} om eksisterende og nye value
    propositions.
  \item \textbf{Vurder det store bildet} -- bransjetrender, konkurrenter,
    nye aktører.
  \item \textbf{Vurder dynamiske effekter} -- hvordan vil kunder,
    konkurrenter og andre reagere?
  \item \textbf{Tenk utover ja/nei-scenarier} -- vær tydelig og nyansert.
  \item \textbf{Vurder hva rammeverk kan og ikke kan} -- bygg bro mellom
    dem heller enn å tvinge inn én linse.
\end{enumerate}

\subsection{Eksamensrelevans}
\begin{itemize}
  \item Demonstrere helhetsforståelse ved å koble forelesningene sammen.
  \item Bruke flere rammeverk komplementært (ikke bare ett).
  \item Vise metarefleksjon: hva kan teorien, hva kan den ikke.
  \item Underbygge alle påstander om virksomheten med kildekritisk evidens.
  \item Lande på en \emph{begrunnet} anbefaling, ikke en lærebokoppsummering.
\end{itemize}
